\begin{examplebox}
	\textbf{\textcolor{CExample}{例 1（基础）}}

	\textbf{\textcolor{CExample}{题目：}}
	已知$f(x)=x^2-4x+5$，$x\in[0,3]$。若$\forall x\in[0,3]$，$a\ge f(x)$恒成立，求实数$a$的最小值。

	\textbf{\textcolor{CExample}{解题步骤：}}
	\begin{description}[style=nextline, leftmargin=2.8em, labelsep=0.5em, itemsep=0.45em]
		\item[\textbf{第一步（识别结构）：}] 不等式为$a\ge f(x)$，直接适用结论，需求$f(x)$在$[0,3]$上的最大值。
			\par\textbf{理由：} 由结论，$a\ge f(x)$恒成立$\iff a\ge \max_{x\in[0,3]} f(x)$。
		\item[\textbf{第二步（求最大值）：}] $f(x)=x^2-4x+5=(x-2)^2+1$，对称轴$x=2$。计算$f(0)=5$，$f(2)=1$，$f(3)=2$，故$\max f(x)=5$。
			\par\textbf{理由：} 二次函数在闭区间上比较端点和顶点可得最大值。
		\item[\textbf{第三步（得范围）：}] 由等价关系，$a\ge5$，故$a$的最小值为$5$。
			\par\textbf{理由：} 代入结论即得参数范围。
	\end{description}

	\textbf{\textcolor{CExample}{关键结论：}}
	\textbf{$a$的最小值为$5$。}

	\medskip
	\textbf{\textcolor{CExample}{例 2（稍变形）}}

	\textbf{\textcolor{CExample}{题目：}}
	若$\forall x\in[1,3]$，$a\le x+\dfrac{4}{x}$恒成立，求实数$a$的最大值。

	\textbf{\textcolor{CExample}{解题步骤：}}
	\begin{description}[style=nextline, leftmargin=2.8em, labelsep=0.5em, itemsep=0.45em]
		\item[\textbf{第一步（识别结构）：}] 不等式为$a\le g(x)=x+\frac{4}{x}$，适用结论，需求$g(x)$在$[1,3]$上的最小值。
			\par\textbf{理由：} $a\le g(x)$恒成立$\iff a\le \min_{x\in[1,3]} g(x)$。
		\item[\textbf{第二步（求最小值）：}] $g'(x)=1-\frac{4}{x^2}=\frac{x^2-4}{x^2}$，令$g'(x)=0$得$x=2$（$x=-2$舍）。计算$g(1)=5$，$g(2)=4$，$g(3)=3+\frac{4}{3}\approx 4.33$，故$\min g(x)=4$。
			\par\textbf{理由：} 利用导数找出极值点，比较端点得最小值。
		\item[\textbf{第三步（得范围）：}] 由等价关系，$a\le4$，故$a$的最大值为$4$。
			\par\textbf{理由：} 代入结论得参数范围。
	\end{description}

	\textbf{\textcolor{CExample}{关键结论：}}
	\textbf{$a$的最大值为$4$。}

	\medskip
	\textbf{\textcolor{CExample}{例 3（综合应用）}}

	\textbf{\textcolor{CExample}{题目：}}
	若关于$x$的不等式$ax\le x^3+1$对任意$x\in[1,2]$恒成立，求实数$a$的取值范围。

	\textbf{\textcolor{CExample}{解题步骤：}}
	\begin{description}[style=nextline, leftmargin=2.8em, labelsep=0.5em, itemsep=0.45em]
		\item[\textbf{第一步（分离参数）：}] 当$x\in[1,2]$时$x>0$，两边同除以$x$得$a\le x^2+\dfrac{1}{x}$。记$g(x)=x^2+\frac{1}{x}$，问题化为$a\le g(x)$恒成立。
			\par\textbf{理由：} 分离参数，注意不等号方向不变。
		\item[\textbf{第二步（求最小值）：}] $g'(x)=2x-\frac{1}{x^2}=\frac{2x^3-1}{x^2}$。令$g'(x)=0$得$x=\sqrt[3]{\frac12}\approx0.79\notin[1,2]$，故$g'(x)>0$恒成立，$g(x)$在$[1,2]$上单调递增。最小值在$x=1$处取得：$g(1)=1+1=2$。
			\par\textbf{理由：} 导数判断单调性确定最小值点。
		\item[\textbf{第三步（得范围）：}] 由等价关系，$a\le2$，故$a$的取值范围是$(-\infty,2]$。
			\par\textbf{理由：} 代入结论。
	\end{description}

	\textbf{\textcolor{CExample}{关键结论：}}
	\textbf{$a\le2$。}
\end{examplebox}
